\documentclass{turabian}
\usepackage{hyperref}
\hypersetup{
    colorlinks=true,
    linkcolor=blue,
    filecolor=magenta,      
    urlcolor=blue,
}
% This template is no longer required to utilize the Turabian LaTeX class.  It is provided as a guide for demonstrating the class's capabilities.  For the latest version of this template and the Turabian class, visit the templates gallery at Overleaf.com.


\begin{document}


% ---------- TITLE PAGE -----------

\centering{
\vspace{30mm}
\textbf{{\huge Computer Science Portfolio Report}\\*
\vspace{30mm}
\\*
\vspace{6mm}
\\*
\vspace{6mm}
{\huge Muhammad Shaheer}\\*
\vspace{6mm}
{\large GH1042020}\\*
\vspace{110mm}
\\*
{\large Job Position being applied for :\href{https://www.stepstone.de/stellenangebote--Praktikant-Werkstudent-IT-Admin-m-w-d-Berlin-torq-partners-People-GmbH--14147529-inline.html}{ IT Admin}}\\*
\vspace{6mm}
{\large Bsc Computer Science student}\\*
\vspace{3mm}
{\large Gisma University Of Applied Sciences, Berlin, Germany}\\*
\vspace{3mm}
{\large realshaheerkhan@gmail.com   |  +49 2631 8308757}\\*
\vspace{3mm}
https://muhammadshaheer2006.github.io/My-Portfolio.github.io/\\*
\vspace{3mm}
https://github.com/MuhammadShaheer2006?tab=repositories\\*}
        }

\pagestyle{empty}
\pagebreak


% ------- TABLE OF CONTENTS -------

\papercontents

Introduction\dotfill 3\\*
Portfolio URLs\dotfill 3\\*
Structure Of Portfolio \dotfill 4\\*
Design Decisions \dotfill 4\\*
CV Design \dotfill 5\\*
GitHub Project Design \dotfill 6\\*
Tools and Technologies \dotfill 8\\*
Reflection, Strengths and Weaknesses \dotfill 9\\*
Conclusion \dotfill 9\\*
\vfill
\pagebreak


% --------- PAPER CONTENT ---------

\content

\section{\large Introduction}
\vspace{-4mm}

This report is accompanied by my computer science portfolio, it is submitted as part of the application for a working student position.
The Pages on my github and portfolio include my CV as well as my own personal projects and a general overview of my technical skills as an upcoming DevOps Engineer.\\*
This report is being written to explain the designs, structure and the content of my portfolio, and explain why I made certain choices when i did, as well as to list the tools I used in making my portfolio and projects. Furthermore, I will also talk about what can be improved and any changes I would like to make in the future.

\section{\large Portfolio URLs}
\vspace{-4mm}

{\Large Portfolio Website}\\*
\vspace{3mm}
My Portfolio website is currently active on GitHub and you can access it here:
https://muhammadshaheer2006.github.io/My-Portfolio.github.io/\\*
\vspace{3mm}
{\Large Repository}\\*
\vspace{3mm}
You can check out my repository for my projects here:
https://github.com/MuhammadShaheer2006?tab=repositories\\*
{\Large CV github}\\*
\vspace{3mm}
You can check out my CV written in LATEX on my github:
https://github.com/MuhammadShaheer2006/Overleaf-CV\\*
As detailed in the repository structure, you can access my 
\href{https://muhammadshaheer2006.github.io/My-Portfolio.github.io/CV%20(1).pdf}{Curriculum Vitae (PDF)} 
directly via my live portfolio website.\\*

\section{\large Structure Of Portfolio}\\*
This Portfolio is divided into 5 major sections namely Home, About, Skills, Projects, Contact:\\*
1- Home: The First page you'll see where a summary of my projects, skills along with my own picture is available, this page has a summary of everything the portfolio has and was made as simple as possible for easy and quick navigation and the style is kept minimalist. My CV is also present on the page in pdf form\\*
2- About: Right next to the home button at the top is the About section where I explain a bit more about myself, my goals as well as my future plans. There is a contact form at the end of the page incase you want to reach out to me. My CV is also on this page in pdf form\\* 
3- Skills: Next to the about tab is the Skills tab, it talks about me own personal skills and everything I know such as HTML/CSS, C++, Python etc. As time goes on I plan on expanding this section to new heights.\\*
4- Projects: Next to the Skills Section is the Projects section showcasing how I used the languages I know from my skills section in real projects, It includes projects like a password generator, a calculator, a guessing game etc.\\*
5- Contact: Last but not least is the contact section, In this section i have my own contact details if you wish to ask me questions or know more details, one some pages I've also listen boxes for you to send me your email so i can contact you too. But generally through this section you can get intouch with me to know more. 

\section{\large Design Decisions}
{\large Multi Page Layout }\\*
For my Portfolio website having a single page layout is good, which is why I made the home page a bit bigger then the other, but at the end of the day a multi page layout was used as if you wanted to know a bit more about me, or my skills projects or contact, you can do so easily.\\*
{\large Minimalist Design }\\*
I decided to make my website a bit more minimalist and straight forward, noone wants to be forced to look at 300 pictures of the employee or read thousands of lines about their entire life story. For this I used a simplistic design thats easy on the eyes and tried to keep most things short and straightforward.\\*
{\large Quick Contact Button }\\*
On the top right corner I put a quick contact button. I added this button to make even easier for employers to get intouch with me. Once the "Lets Get Together" button is clicked, it redirects you to the contact page where all my contact information like my email and phone number are.\\*
{\large GitHub Page Hosting }\\*
GitHub was used to host the portfolio website as it is the professional and industry standard for making personal portfolio as well as websites. Using GitHub made the process simple and straightforward allowing me to easily navigate through my repositories and make my portfolio.\\*
{\large LaTeX CV }\\*
I added a LaTeX CV to the website which can be viewed right there by simply the press of a button, I also have a GitHub repository for my CV which can be viewed here: https://github.com/MuhammadShaheer2006/Overleaf-CV\\*

\section{\large CV Design}
CV Repository link: https://github.com/MuhammadShaheer2006/Overleaf-CV\\*
As detailed in the repository structure, you can access my 
\href{https://muhammadshaheer2006.github.io/My-Portfolio.github.io/CV%20(1).pdf}{Curriculum Vitae (PDF)} 
directly via my live portfolio website.\\*
My CV on GitHub uses a professional design and simple design for easy viewing and clear objectives. On the left I labeled the simpler parts such as the Contact List, which has all of my Contact Details along with my github and linked in that can be clicked to be taken directly to them.\\*
After the contact list is the Skills List, where I have highlighted all of my Skills in different languages such as Python, C++, LaTeX, HTML/CSS, along with my proficiency in them. Right below the coding skills are the operating systems I work with and how well I know them all. And lastly are the languages I know and speak with my expertise level in them.\\*
Below that are the tool and technologies I work with and lastly my nationality with a QR code which can be scanned to be taken to my GitHub.\\*
On the right side I have an About Me which gives a brief intro of who I am and what my ideals are along with my future plans in the DevOps field. It is kept short and straight forward.\\*
In the Education section I have highlighted my education status soo far over the years as well as what subjects I have studied.\\*
The Achievements tab was made small and simple to be clear and concise of my accomplishments\\*
The General Skills section showcases the practical skills I have acquired over the years the positives of working with me and my capabilities.\\*
Lastly I made a certificate Section talking about the official certificates I have received to give official proof of my coding skills\\*

\section{\large GitHub Project Design}
{\Large 1 - City Guesser}\\*
This project was made to be a fun game to pass the time, it features 10 cities and only 1 try, if you fail, you lose and the program tells you the city and then restarts for another round. It was made in C++ with simple code. It was my first attempt which is why its a bit janky.\\*
https://github.com/MuhammadShaheer2006/City-Guesser\\*
{\Large 2 - Password Generator}\\*
This is a simple random password generator for when you need a good password, it uses lower case, upper case and punctuation to make strong unpredictable passwords that can be used anywhere. It was made in python and can be used as a tool for security when making your passwords.\\*
https://github.com/MuhammadShaheer2006/Password-Generator\\*
{\Large 3 - Calculator}\\*
This project is a calculator that can multiply, divide, add and subtract 2 numbers. You can use it easily and input your numbers, then choose whatever operations you want to use. It was also written in C++ in order for me to experiment and learn more about the language.\\*
https://github.com/MuhammadShaheer2006/Calculator\\*
{\Large 4 - Converter}\\*
This is a simple conversion program, you can choose what do you wanna convert, Celsius, Kilograms or Meters. After you've chosen a metric, the program gives the result, for example if you choose celsius, the program will output Kelvin and Fahrenheit temperature. It was also written in Python by me to experiment with different statements.\\*
https://github.com/MuhammadShaheer2006/Conversion-System\\*
{\Large 5 - Task Manager}\\*
This task manager allows users to write down their tasks on the go and add any new tasks they want to do in the day, It is perfect for people with busy schedules. This program was also written in C++ with the purpose of making a to do list for your tasks easier.\\*
https://github.com/MuhammadShaheer2006/Task-Manager\\*
{\Large 6 - Log}\\*
Thus is a time tracking system where you can put your timestamps, for example if your studying and put an alarm for 1 hour, after an hour you can write that down in the log program by telling it what activity did you do and for how long, it'll save that activity and the time given in minutes for you to remember what task you did, when did you log them and how long they lasted.\\*
https://github.com/MuhammadShaheer2006/Log\\*

\section{\large Tools and Technologies}
{\Large HTML/CSS}\\*
These are the core building blocks for the portfolio on the basis of which my entire portfolio website is made, they were used to design every inch of the website along with the the colors and all the features.\\*
{\Large Git and GitHub}\\*
This handles the Version control and remote hosting of all source files. The repository tracks all changes and enables rollback if needed.\\*
{\Large GitHub Pages}\\*
This has allowed me to host my pages online as well as make web hosting simple and accessible. It is whats used for hosting my own portfolio.\\*
{\Large LaTeX}\\*
LaTeX was used to write the CV available on my github, along with this Computer Science Report too. It is has been major help in keeping things professional and clear.\\*
{\Large Visual Studio Code}\\*
This was used to write the programs themselves that were used in the projects, and to fine tune them to make sure my github projects have no errors. This is where they were tested for any issues and fixed.\\*

\section{\large Reflection, Strengths and Weaknesses}
{\Large Strengths}\\*
1 - Multi page Portfolio layout for simple and easy navigation along with more deeper details on each dedicated tab.\\*
2 - Quick contact button for easy contact information in the top right corner.\\*
3 - Completed portfolio website and CV with all necessary details.\\*
4 - LaTeX CV used with industry standard layout\\*
5 - Accessibility with public URLs to easily access the portfolio anywhere.\\*

{\Large Weaknesses and areas of Improvement}\\*
1 - A major weakness of my Report is the lack of more complex projects that really test my coding skills. This will be improved overtime as I learn more and more.\\*
2 - Limited use of java script is also a downside, the website is slightly interactable but this can be further improved upon as I learn.\\*
3 - The Skills section on the portfolio not having more details is also a drawback but that can be changed when im confident enough in my coding abilities to boast about it.\\*
4 - The portfolio repository is also very confusing and it will be improved in the future to make it as simple and user friendly to understand as possible if someone else wants to look at the repository.\\*
5 - Lack of previous working experience can be made up for as I take on new job and better opportunities in my digital coding journey.

\section{\large Conclusion}
In conclusion this report represents how I designed and made my portfolio, CV and the projects in my portfolio and provides their github repository links for viewing them. It represents my current standing as a computer science student who wants to excel in DevOps in this ever evolving digital world. In the future I plan on bettering my current Skills and trying new languages to help me excel in achieving my goals.



{\large LINK TO MY PORTFOLIO WEBSITE :}

https://muhammadshaheer2006.github.io/My-Portfolio.github.io/

{\large LINK TO MY GITHUB REPOSITORIES :}

https://github.com/MuhammadShaheer2006?tab=repositories

Here you can access my portfolio and repository directly.



\pagebreak

\end{document}